\documentclass[11pt,a4paper]{article}
\usepackage[T1]{fontenc}
\usepackage[utf8]{inputenc}
\usepackage{lmodern}
\usepackage[margin=2.5cm]{geometry}
\usepackage{setspace}
\usepackage{parskip}
\setlength{\parskip}{0.7em}
\pagestyle{empty}

\begin{document}

\begin{center}
{\large\textbf{Statement of Interest}}\\[0.4em]
{\normalsize Wealth of Nations at 250 --- Rapporteur Application}\\[0.2em]
{\small Aristeidis Grivokostopoulos \quad|\quad a.grivokostopoulos@lse.ac.uk}\\[0.2em]
{\small Contributor preferences: (1) Besley \& Ghatak \quad (2) Rodrik \quad (3) Voth}
\end{center}

\bigskip

I am a PhD candidate in Economic History at the London School of Economics, funded by the ESRC. My dissertation examines sovereign debt and political institutions across the first age of globalisation, combining archival research with computational text analysis and panel econometrics.

My preferred assignments are Besley and Ghatak, Rodrik, and Voth, and the fit in each case is direct. My paper ``Between Civilisation and Constitution: Greek Parliamentary Debates on Gold Standard Adoption, 1875--1885'', prepared for the Mary Morgan Retirement Workshop (June 2026), applies Morgan's models as mediators framework to show how the gold standard operated in nineteenth-century Greek politics: not as a self-enforcing commitment mechanism but as a model whose authority depended on narrative embedding as much as technical demonstration. This puts me in close engagement with the questions Besley and Ghatak raise about the interplay between values, narrative, and economic institutions: their chapter on moral sentiments in market economies addresses exactly the mechanisms I reconstruct from the Greek parliamentary record. A companion paper on Greek legislative behaviour (1875--1905) uses 755 roll call votes to test whether fiscal shocks shift legislators from clientelist to programmatic behaviour, which is the question Besley and Ghatak address from the theoretical side: when do values and motivated agents reshape economic institutions rather than adapt to them?

Rodrik's chapter raises the question of contextual versus universal application of economic ideas, which is the core of my paper's misapplication claim: the critics of Greece's gold standard adoption were technically correct, and they lost not because they were wrong but because they lacked a counter-model with comparable narrative authority. In Rodrik's terms, the critics were making a proto-mercantilist case, prioritising contextual fit over universal principle, without a framework capable of matching the gold standard's civilisational authority.

Voth's chapter on conflict, colonialism, and the wealth of nations connects to my research on debt restructuring under the Ottoman Public Debt Administration, the Egyptian Caisse de la Dette, and the Greek International Financial Commission, where I study the fiscal infrastructure question from the borrower side using bond-level London Stock Exchange data. All three supervisory regimes followed military defeat and forced fiscal restructuring; my difference-in-differences estimates show significant yield spread convergence in each case, consistent with Voth's argument that conflict-driven institutional change reshaped credit conditions for generations.

As a teaching assistant for EC311 (History of Economic Thought) at LSE, I regularly explain how economic ideas develop, travel, and are contested across time, which is the same task a rapporteur proceedings chapter performs for a conference contribution. I have presented my own research at WEHC (Lund, 2025) and the WEAI Annual Conference (San Francisco, 2025), and I write for audiences across economic history, political economy, and the history of economic thought.

Having active papers in each of the areas these chapters address, I would engage with the contributions as a researcher, not only as a note-taker. I attach my CV and a letter of support from my supervisor, Professor Albrecht Ritschl.

\end{document}
